% !TEX program = xelatex
\documentclass[12pt,letterpaper]{article}

\usepackage[margin=1in]{geometry}
\usepackage{fontspec}
\usepackage{setspace}

\setmainfont{Times New Roman}
\setlength{\parindent}{0.5in}

\begin{document}
\doublespacing

\begin{center}
    \textbf{Personal Statement for the University of Washington}\\
    \textbf{Master in Social Work Program:}\\
    \textit{On the Topic of Free Will}\\
    Nathan Wesley
\end{center}

It is the Fall of 2024. You are 25 years old, and your particularly stubborn yet philosophical roommate asks “Does free will exist?” You say, “Yes, of course free will exists, but also no.” The two of you engage in an unwelcome hour-long dialogue. This conversation is frustrating because you believe strongly in your opinion but lack the vocabulary to articulate it fully. Two years later, you have a better answer, and thankfully the two of you are no longer roommates.

We must first define ‘free will’. Let us begin with the latter half. It may be easier to describe ‘will’ through its negative. ‘Will’ is not instinct, reflex, or blind obedience. ‘Will’ is neither random nor coincidental. ‘Will’ is the conscious intent to act. This shall be an appropriate definition of ‘will’.

On to the former. For this argument, there are two relevant prepositional complements to the word ‘free.’ Firstly, one may be ‘free to’ do something, meaning the individual has both the permission (without external barriers) and the genuine capability (internal agency and resources) to choose and carry out that action. I find that it is easy to equate ‘free will’ with ‘free to’. Every, day we make many conscious choices that we are ‘free to’ make. Proponents of “yes, free will exists” ignore the second preposition: ‘free from.’ ‘Free from’ means being completely without or unaffected by something. We are never truly ‘free from’ any condition.

It is now months after the unpleasant conversation with your roommate. Opaque, winter gray paints the Seattle sky. The sun is gone, and everything is wet. Though, this doesn’t bother you. Through intense meditation, you’ve unlocked the ability to “really see colors”. You’ve successfully charted out a complete map of human consciousness on the walls of your rental unit. Words in your mind arrive in storms. You’re speed-reading the dictionary. You haven’t slept in two days. You step into your therapist’s office and proclaim, “I’m cured!” Your therapist diagnoses you with Bipolar Disorder.

Receiving a bipolar diagnosis is a particularly destabilizing experience (aside from, actually, experiencing bipolar disorder). Life becomes wholly paradoxical. The past is clearer now. A real explanation for years of sleepless nights and days full of sleep. You’re no longer lazy or crazy, just sick. Life makes sense now. Your future becomes uncertain, and that is scary. Denial arrives shortly after.

It is quite easy to lose your sense of self after a diagnosis. You reflect on all the decisions that you make and wonder if that was really you or the “bipolar you.” After spending an ungodly amount of hours researching the topic, you realize that there is no cure for Bipolar Disorder. You take your medications for some time, feel better, and believe you’re cured. You relapse into an episode and repeat this cycle for months and months. Some time later you realize that you will probably never be ‘free from’ this disease.

Months later, you are comfortably in a nihilistic arc of existence. Seeking a concrete meaning for life, you sign up to volunteer at an emergency shelter for young adults. You’re not sure where exactly the idea came from. Maybe someone mentioned the topic in passing by coincidence, or perhaps a random act of kindness made you feel particularly altruistic that day. Nevertheless you choose to volunteer your evenings there.

Homelessness is a controversial topic. You discover that even the term “homeless” is controversial. There are so many reasons why someone can experience housing insecurity. It might be hard for somebody to hold a job due to medical complications, while someone else may be on the streets due to familial rejection of their gender identity. A lot of people are fleeing from scary situations. You find out about so many possible systemic inequalities that can result in homelessness, reasons that are hard to be ‘free from.’ This can be disparaging.

Yet, you find so many more reasons to believe in people. There is a perpetual kindness that permeates the space. People choose to listen to and support each other. They choose to tolerate and accept one another. They choose to be present for themselves. It’s hard work. Every evening you’re there, you witness people consciously choose to do something good. It seems that waves of life can run you ashore to unexpected places. But even when you’re lost, there exists an internal motivation to act. Free will exists sometimes in some places, it seems.

Years later, you’ve been stable for quite some time now, and you don’t feel sick. You feel like a person. You have graduated from the school of nihilism (you are an absurdist now). You’re still not sure if there exists a meaning for life, or what your Dharma is. All you know is that you find the process of therapy incredibly interesting. The world has a lot of problems, and you think you can help somehow. You tell your therapist that you want to sit in her chair. So you write this.

\end{document}
